\documentclass[11pt]{article}
\usepackage[margin=1in]{geometry}
\usepackage{amsmath,amssymb,amsthm,booktabs}
\usepackage[hidelinks]{hyperref}
\newcommand{\F}{\mathbb{F}_2}
\newcommand{\Z}{\mathbb{Z}_4}

\title{Orbit unions and exchange walks:\\ seven new lower bounds for binary codes}
\author{Zhiyang Chen \and Ligeng Zhu \and Fan Long}
\date{October 2026}

\begin{document}
\maketitle

\begin{abstract}
Let $A(n,d)$ be the largest size of a binary code of length $n$ and minimum distance $d$. We give explicit codes
that show $A(17,6)\ge260$, $A(21,6)\ge3072$, $A(24,10)\ge144$, $A(25,10)\ge256$, $A(27,6)\ge81920$,
$A(28,6)\ge139264$ and $A(28,12)\ge192$. Six of the codes are unions of orbits of a group of isometries: we choose
the group, build the graph of orbits, and select orbits by a maximum independent set search. Most of the groups
are translation groups of $\Z$-linear codes under the Gray map. The seventh code was found by an exact exchange walk
that starts from such a code. Each code is given as a file, and its minimum distance is checked directly.
\end{abstract}

\section{Results}

Table~\ref{tab:results} lists the bounds; ``old'' is the best value known to us before this work. Four of the bounds
were sent to A.~E.~Brouwer on 2026-10-21 and are now in his table~\cite{Br}. Since $A(n-1,d)\ge A(n,d)/2$
(shortening), we also get $A(26,6)\ge40960$ and $A(25,6)\ge20480$ (old $32768$ and $17920$).

\begin{table}[ht]
\centering
\begin{tabular}{rrrl}
\toprule
$(n,d)$ & old & new & construction \\
\midrule
$(17,6)$ & 258 \cite{Mil} & \textbf{260} & exchange walk from an orbit union \\
$(21,6)$ & 2560 & \textbf{3072} & 48 cosets of a $\Z$-code of order 64, with 4 extra coordinates \\
$(24,10)$ & 136 \cite{LO} & \textbf{144} & 18 cosets of a 3-dimensional subcode of the Golay code \\
$(25,10)$ & 196 \cite{OOH} & \textbf{256} & 16 cosets of a $\Z$-code of order 16, with 3 extra coordinates \\
$(27,6)$ & 65536 & \textbf{81920} & 20 cosets of a $\Z$-code of order 4096, with 3 extra coordinates \\
$(28,6)$ & 131072 & \textbf{139264} & 34 cosets of the same $\Z$-code, with 4 extra coordinates \\
$(28,12)$ & 178 & \textbf{192} & 12 cosets of a 4-dimensional binary linear code \\
\bottomrule
\end{tabular}
\caption{New lower bounds. The bounds for $(21,6)$, $(27,6)$, $(28,6)$ and $(28,12)$ are already in~\cite{Br}.}
\label{tab:results}
\end{table}

\section{The method}

\paragraph{Orbit unions.} Let $\Gamma$ be a group of isometries of $(\F^n,d_H)$. A union of $\Gamma$-orbits
$O_1\cup\dots\cup O_m$ has minimum distance at least $d$ if and only if each $O_i$ has internal minimum distance at
least $d$ and $d(O_i,O_j)\ge d$ for $i\ne j$. Since $\Gamma$ acts by isometries, $d(O,O')=\min_{y\in O'}d(x,y)$ for
any fixed $x\in O$. So the largest $\Gamma$-invariant code is a maximum-weight independent set in the graph whose
vertices are the orbits with internal distance $\ge d$ (weight $|O|$), two orbits being adjacent when their distance is
$<d$. We choose $\Gamma$ so that this graph is small, and search it with tabu search, the CHILS heuristic~\cite{CHILS}
or an exact clique solver. We use three kinds of groups:
\begin{itemize}
\item \emph{Translations by a binary linear code $L$.} The orbits are the cosets $x+L$; the internal distance is the
minimum weight of $L$, and $d(x+L,y+L)$ is the minimum weight of $x+y+L$.
\item \emph{Translations by a $\Z$-linear code $K\le\Z^m$.} The Gray map $\phi:\Z\to\F^2$, $0\mapsto00$, $1\mapsto01$,
$2\mapsto11$, $3\mapsto10$, is an isometry from the Lee metric to the Hamming metric~\cite{HKCSS}. So the maps
$\phi(x)\mapsto\phi(x+k)$, $k\in K$, are isometries of $\F^{2m}$. We also add $e$ binary coordinates, i.e.\ we work in
$\F^e\times\Z^m$ with the distance $d_H+d_{\mathrm{Lee}}$, and take unions of cosets of $\{0\}\times K$. A coset has
internal distance equal to the minimum Lee weight of $K$.
\item \emph{Coordinate permutations}, e.g.\ a cyclic group generated by one permutation.
\end{itemize}
A useful choice for $K$ is a subgroup of small index in a $\Z$-linear code of large minimum Lee distance: its cosets
inside that code are pairwise compatible, and the extra coordinates and the cosets outside the code give room for more.

\paragraph{Exact exchange walks.} Let $C$ be a code with minimum distance $d$. A move removes a set $R$ of $k\le60$
codewords near a random point, and computes (exactly, with a node limit) a largest set $R'$ of words that are at distance
$\ge d$ from $C\setminus R$ and from each other. If $|R'|\ge|R|$ the move is accepted. So the walk moves on the plateau
of codes of equal size until a move gives $|R'|>|R|$. Codes from one large group are usually rigid (no equal-size move
exists), so we start walks from many structurally different codes of the best known size or slightly below it, and every
code that a walk grows becomes a new start. In practice we puncture a random coordinate, walk on the $(n-1,M,d-1)$
code, and add the parity bit back; a node limit makes a small fraction of the searches inexact. This affects only the
search: the final code is checked directly.

\section{The codes}

For even $d$, an overall parity bit turns an $(n-1,M,d-1)$ code into an $(n,M,d)$ code (``parity extension'').
In the files of the $\Z$ constructions the $e$ extra binary coordinates come first, followed by the Gray images;
the only exception is $A(21,6)$, where the $16$ Gray coordinates come first, then the $4$ extra coordinates and the
parity bit. The coset lists and the formats are in the data files.

\paragraph{$A(17,6)\ge260$.} Let $\sigma=(0\,1\,2\,3\,4)(5\,6\,7\,8\,9)(10\,11\,12\,13\,14)$ act on $16$ coordinates.
CHILS on the graph of $\langle\sigma\rangle$-orbits gave a $\langle\sigma\rangle$-invariant $(16,257,5)$ code; its parity
extension is a $(17,257,6)$ code. An exchange walk on it reached a $(17,259,6)$ code within a walk of $1.3$ million moves,
and a second walk from that code reached a $(17,260,6)$ code after $0.77$ million moves. The final code has no symmetry that we
imposed; its kernel $\{v: v+C=C\}$ is trivial.

\paragraph{$A(21,6)\ge3072$.} Let $O\le\Z^8$ be the octacode ($|O|=256$, minimum Lee distance $6$; its Gray image is the
Nordstrom--Robinson code~\cite{HKCSS}). For a subgroup $H\le O$ of order $64$, a union of $48$ cosets of $\{0\}\times H$
in $\F^4\times\Z^8$ is a $(20,3072,5)$ code, and its parity extension is a $(21,3072,6)$ code. (The best union of
cosets of the full octacode in this setting has only $2560$ words.)

\paragraph{$A(24,10)\ge144$.} Let $L$ be a $3$-dimensional subcode of the binary Golay code $[23,12,7]$ with minimum
weight $11$. A union of $18$ cosets of $L$ in $\F^{23}$ is a $(23,144,9)$ code (the cosets need not lie in the Golay code; here
none does); its
parity extension is a $(24,144,10)$ code. It is not linear, since $144$ is not a power of $2$.

\paragraph{$A(25,10)\ge256$.} Let $C\le\Z^{11}$ be generated by
\[
g_1=(2,2,2,0,2,0,0,0,2,0,0),\quad g_2=(2,3,3,1,0,0,1,2,1,1,0),\quad g_3=(1,0,2,1,2,1,3,3,0,0,1).
\]
Then $|C|=32$ and $C=\langle g_1\rangle\oplus\langle g_2\rangle\oplus\langle g_3\rangle$ has minimum Lee distance $10$.
Let $K=\{ag_1+bg_2+cg_3: a+b+c \text{ even}\}$, of order $16$ and minimum Lee distance $10$. A union of $16$ cosets of
$\{0\}\times K$ in $\F^3\times\Z^{11}$ is a $(25,256,10)$ code.

\paragraph{$A(27,6)\ge81920$ and $A(28,6)\ge139264$.} Let $C=\ker\Phi\le\Z^{12}$, where
$\Phi:\Z^{12}\to\mathbb{Z}_2^2\times\Z^4$ is the homomorphism given in the data file \texttt{z12\_1.h}; then $|C|=2^{14}$ and
$C$ has minimum Lee distance $6$. Let $\psi:\Z^{12}\to\mathbb{Z}_2^2$, $\psi(x)=\sum_{i:\,x_i\text{ odd}}\psi_i$, with
$(\psi_1,\dots,\psi_{12})=(1,2,1,1,2,1,3,0,0,0,2,1)$ (an integer $t\in\{0,1,2,3\}$ stands for its binary digits). Then
$C'=C\cap\ker\psi$ has order $4096$. A union of $20$ cosets of $\{0\}\times C'$ in $\F^3\times\Z^{12}$ is a
$(27,81920,6)$ code, and a union of $34$ cosets in $\F^4\times\Z^{12}$ is a $(28,139264,6)$ code. Both rebuild exactly from
$(\Phi,\psi)$ and the coset lists in the data files.

\paragraph{$A(28,12)\ge192$.} The code is a union of $12$ cosets of a $4$-dimensional linear code $L\le\F^{28}$ of
minimum weight $12$, and it lies in a coset of a $12$-dimensional linear code. It was found by a coset search with a
smaller kernel taken inside a good linear code of length $27$, followed by parity extension.

\section{Verification and data}

Each code is a list of binary words. Its size, its length and the absence of two words at distance $<d$ were checked by
at least two independent programs; this finite check is the proof. The structure stated above was also checked:
the $\Z$ codes were rebuilt from their parameters, and the cosets and kernels were recomputed from the files.

\begin{center}
\begin{tabular}{ll}
\toprule
file & SHA-256 (first 16 hex digits) \\
\midrule
\texttt{A\_17\_6\_260.txt} & \texttt{489b6f11abdc0ae7} \\
\texttt{A\_21\_6\_3072.txt} & \texttt{dfc23bec786e9375} \\
\texttt{A\_24\_10\_144.txt} & \texttt{d723cfa082ea9230} \\
\texttt{A\_25\_10\_256.txt} & \texttt{8c64ab1860a01dd8} \\
\texttt{A\_27\_6\_81920.txt} & \texttt{0038594ad6dd0b20} \\
\texttt{A\_28\_6\_139264.txt} & \texttt{df5c845a6c748619} \\
\texttt{A\_28\_12\_192.txt} & \texttt{5149782b3253fe49} \\
\bottomrule
\end{tabular}
\end{center}

The code files, the group and coset data, and the search programs are in the data record of this paper (\url{https://doi.org/10.5281/zenodo.23599650}).

\emph{Use of AI.} AI tools were used in the search and in the writing of this paper. The correctness of the
bounds does not depend on them: each code is checked directly, as described above, and the authors guarantee the
stated bounds.

\begin{thebibliography}{9}
\bibitem{Br} A.~E.~Brouwer, Table of general binary codes, \url{https://aeb.win.tue.nl/codes/binary-1.html}.
\bibitem{CHILS} E.~Gro{\ss}mann, K.~Langedal, C.~Schulz, Accelerating reductions using graph neural networks and a new
concurrent local search for the maximum weight independent set problem, arXiv:2412.14198.
\bibitem{HKCSS} A.~R.~Hammons, P.~V.~Kumar, A.~R.~Calderbank, N.~J.~A.~Sloane, P.~Sol\'e, The $\mathbb{Z}_4$-linearity of
Kerdock, Preparata, Goethals, and related codes, \emph{IEEE Trans.\ Inform.\ Theory} 40 (1994) 301--319.
\bibitem{LO} A.~Laaksonen, P.~R.~J.~\"Osterg\aa rd, Constructing error-correcting binary codes using transitive
permutation groups, \emph{Discrete Appl.\ Math.} 233 (2017) 65--70.
\bibitem{Mil} M.~Milshtein, A new two-error-correcting binary code of length 16, \emph{Cryptogr.\ Commun.} (2019).
\bibitem{OOH} K.~Ota, T.~Osa, T.~Harada, Three improved lower bounds for binary error-correcting codes,
arXiv:2610.11609 (2026).
\end{thebibliography}
\end{document}
